\documentclass[11pt]{article}
\usepackage[margin=0.8in]{geometry}
\usepackage{booktabs,hyperref}
\hypersetup{colorlinks=true,urlcolor=blue}
\emergencystretch=3em
\title{Impala performance experiment: Step 3}
\author{Float64 GPU material-history prototype}
\date{September 29, 2026}
\begin{document}\maketitle
\section*{Decision}
GPU benefit depends on batch size. It is substantially slower than the compiled
CPU at small batches, modestly faster at 50 parameter sets, and substantially
faster for large batches. Keep this as an experimental benchmark adapter rather
than a default GPU backend. In the matched seed-42 calibration, GPU sampling
was 13.5\% slower than the compiled CPU. The full 50-temperature pooled study passed
the predeclared posterior-quality gates. The known hierarchical convergence
issue remains unresolved; short hierarchical runs certify execution only.
\section*{Implementation and setup}
The prototype assigns one CUDA thread to each material history and advances all
100 steps in that thread. It uses double precision, disables fused multiply-add,
and does not enable fast math. Cu PTW equations match the CPU implementation;
GPU transcendental functions can still round differently. Interpolation and
MCMC stay on the CPU. Conditions, parameters and initial states are transferred
every call; full histories and final state are returned. Timings include host
preparation, allocation, transfers and interpolation. No persistent device grid
cache is used; CuPy can reuse memory-pool allocations.

Hardware: NVIDIA RTX PRO 5000 Blackwell Generation Laptop GPU, compute capability
12.0, about 24 GB memory. CuPy 14.2.0 and CUDA runtime/NVRTC 13.2.51 were installed
as optional virtual-environment packages. The system toolkit was root-readable
only, so no system permission changes were made. An existing GPU process was
left untouched; sequential timings are observations under this machine's load.
\section*{Numerical accuracy}
All 91 tests passed in 113.57 seconds, including eight explicit GPU hardware tests.
Tests cover pooled/hierarchical batches, full histories and final fields,
continued states, edited inputs, different integration grids, small-p branches,
strain-rate thresholds and invalid PTW parameters. The original 256-case fixture
had maximum prediction error 5.97e-16 Mbar and maximum weighted log-likelihood
error $1.75\times10^{-10}$. Prediction tolerances are relative $10^{-8}$
and absolute $10^{-10}$ Mbar. The log likelihood absolute tolerance is $10^{-6}$.
Baseline fixtures were not regenerated.

Unlike step 2, exact GPU chain equality is not required or expected. Tiny
rounding changes can eventually change MCMC accept/reject decisions. We assess
posterior agreement and effective samples per second in addition to numerical
fixtures, using criteria declared before calibration runs.
\section*{Transfer-inclusive model timing}
Thirty paired, alternating CPU/GPU evaluations per batch, after warm-up.
CPU is the compiled step-2 implementation. Times are milliseconds; speedup is
CPU time divided by GPU time. Kernel-only times were measured separately with
CUDA events and exclude host work and transfers.
\begin{center}\begin{tabular}{rrrrr}\toprule
Batch & CPU (ms) & GPU (ms) & Kernel (ms) & Speedup\\\midrule
1 & 0.446 & 2.344 & 1.811 & 0.19\\
4 & 0.643 & 2.293 & 1.912 & 0.28\\
8 & 0.898 & 2.345 & 1.915 & 0.38\\
50 & 3.606 & 3.099 & 2.115 & 1.16\\
200 & 13.104 & 4.512 & 2.136 & 2.90\\
1000 & 67.609 & 13.887 & 2.121 & 4.87\\
\bottomrule\end{tabular}\end{center}
A batch of $B$ parameter sets expands to $6B$ copper experiment histories.
Near-constant kernel latency at these sizes explains why small batches lose.
Startup/compilation is excluded from these warmed timings. First-call metadata
is recorded, but the CUDA disk cache may already be populated.
An instrumented 3,000-state pooled CPU run at 50 temperatures recorded
3,288 model calls, with mean batch size 39.72 and range 1--50. Constraint
filtering therefore makes real calls smaller than the fixed batch-50 test.
This early-chain sample helps explain the crossover; it is not a full-run profile.
\section*{Calibration timing}
Small-batch smoke runs: seed 42, 3,000 states, eight temperatures.
\begin{center}\begin{tabular}{lrrr}\toprule
Mode & CPU (s) & GPU (s) & Speedup\\\midrule
pooled & 3.55 & 8.90 & 0.40\\
hier & 8.01 & 12.32 & 0.65\\
\bottomrule\end{tabular}\end{center}
The larger study used 40,000 states and 50 temperatures. Fresh compiled CPU
seed 42 took 186.68 seconds; GPU seed 42 took 211.88 seconds,
a speedup of 0.88x. This single current CPU timing is not a multi-seed
throughput confidence interval. GPU sampling times:
\begin{center}\begin{tabular}{rr}\toprule
Seed & Sampling seconds\\\midrule
42 & 211.88\\
123 & 215.08\\
456 & 211.44\\
789 & 202.79\\
\bottomrule\end{tabular}\end{center}
The fresh compiled CPU control was also compared with the saved original seed-42
chain. Exact agreement for the checked arrays: True.
\section*{Pooled posterior quality}
Four independent GPU seeds (42, 123, 456, 789) are compared with the preserved
original CPU study at the same configuration. Discard the first 20,000 states
of each chain; analyze all remaining states without thinning. Diagnostics cover
10 material parameters and six noise standard deviations.

Predeclared gates: rank R-hat below 1.01; bulk and tail ESS at least 400;
MCSE(mean)/SD below 0.05; posterior-mean differences at most four combined MCSEs
and 0.1 pooled posterior SD; maximum prediction-mean difference at most 0.1 noise
SD. Prediction checks use the same 200 retained draws per run as the baseline.
\begin{center}\begin{tabular}{lrr}\toprule
Diagnostic & Original CPU & GPU\\\midrule
Maximum R-hat & 1.003 & 1.003\\
Minimum bulk ESS & 2253 & 2253\\
Minimum tail ESS & 3345 & 3345\\
Minimum bulk ESS/s & 2.342 & 2.679\\
\bottomrule\end{tabular}\end{center}
Maximum posterior mean difference: $4.21\times10^{-13}$ combined MCSEs
and $2.12\times10^{-15}$ posterior SD. Maximum prediction-mean difference:
$2.47\times10^{-15}$ noise SD. All quality gates passed. Credible-interval
endpoint differences are recorded separately in the comparison CSV.
All four parameter trajectories and acceptance counters matched exactly.
Noise-variance differences were below $3.7\times10^{-18}$; saved-prediction
differences were below $7.1\times10^{-17}$ Mbar. Thus the combined parameter
and noise arrays are not bitwise identical despite identical parameter draws.

The ESS/s table uses total sampling seconds across four runs. Original CPU
runtimes are historical and precede both CPU optimizations. Thus the observed
GPU/original ESS/s ratio (1.14) must not be described as the gain over
the current compiled CPU. Only the fresh seed-42 timing provides that direct
runtime comparison. The pooled result does not establish hierarchical quality
or generalize to other material models.
\section*{Tracking and reproduction}
Prototype commit: \texttt{7817d9a}. Branch: \texttt{perf/gpu-material-history}.\\
Final checkpoint: \texttt{perf-step3-20260929}. All commits and tags are local.
Production CPU library code and defaults are unchanged. Source, tests, pinned
optional requirements, declared gates, compact results, CSV diagnostics and
LaTeX source are in Git. Raw traces, logs and this PDF are in workspace
\texttt{benchmark/changes/step3/}, with a versioned hash manifest.

See \path{docs/performance/STEP3.md} for individual commands and limitations.
The complete sequential calibration order is reproducible with
\path{benchmarks/gpu_suite.py}; quality analysis uses
\path{benchmarks/gpu_quality.py} in the ArviZ environment. GPU tests require
\texttt{IMPALA\_TEST\_GPU=1}; normal CPU CI skips them. Installation reference:
\url{https://docs.cupy.dev/en/stable/install.html}.
\end{document}
